% TODO make hw stylesheet
\documentclass[10pt]{scrartcl}
\usepackage[a4paper, total={6in, 8in}]{geometry}
\usepackage[usenames,dvipsnames,svgnames]{xcolor}
\usepackage[shortlabels]{enumitem}
\usepackage[framemethod=TikZ]{mdframed}
\usepackage{amsmath,amssymb,amsthm}
\usepackage[colorlinks]{hyperref}
\usepackage{multicol}
\usepackage{tabularx}

\hypersetup{
	colorlinks=true,
	linkcolor=blue,
	filecolor=magenta,      
	urlcolor=magenta,
	pdftitle={MAT 362 HW}
}

\usepackage{microtype}
\usepackage{mathtools}
\usepackage[headsepline]{scrlayer-scrpage}
\usepackage{thmtools}
\usepackage[textsize=footnotesize]{todonotes}

\mdfdefinestyle{mdbluebox}{roundcorner=10pt,innerbottommargin=9pt,
	linecolor=blue,backgroundcolor=TealBlue!5,}
\declaretheoremstyle[headfont=\sffamily\bfseries\color{MidnightBlue},
mdframed={style=mdbluebox},]{thmbluebox}

\mdfdefinestyle{mdbloobox}{frametitlefont=\bfseries,innerbottommargin=8pt,
	nobreak=true,backgroundcolor=TealBlue!5,linecolor=blue,}
\declaretheoremstyle[headfont=\bfseries\color{MidnightBlue},
mdframed={style=mdbloobox},headpunct={\\[3pt]},postheadspace=0pt,]{thmbloobox}

\mdfdefinestyle{mdredbox}{frametitlefont=\bfseries,innerbottommargin=8pt,
	nobreak=true,backgroundcolor=Salmon!5,linecolor=RawSienna,}
\declaretheoremstyle[headfont=\bfseries\color{RawSienna},
mdframed={style=mdredbox},headpunct={\\[3pt]},postheadspace=0pt,]{thmredbox}

\mdfdefinestyle{mdgreenbox}{linecolor=ForestGreen,backgroundcolor=ForestGreen!5,
	linewidth=2pt,rightline=false,leftline=true,topline=false,bottomline=false,}
\declaretheoremstyle[headfont=\bfseries\sffamily\color{ForestGreen!70!black},
mdframed={style=mdgreenbox},headpunct={ --- },]{thmgreenbox}

\mdfdefinestyle{mdorangebox}{linecolor=RedOrange,backgroundcolor=RedOrange!5,
	linewidth=2pt,rightline=false,leftline=true,topline=false,bottomline=false,}
\declaretheoremstyle[headfont=\bfseries\sffamily\color{RedOrange!70!black},
mdframed={style=mdorangebox},headpunct={ --- },]{thmorangebox}

\mdfdefinestyle{mdblackbox}{linecolor=black,backgroundcolor=RedViolet!5!gray!5,
	linewidth=3pt,rightline=false,leftline=true,topline=false,bottomline=false,}
\declaretheoremstyle[mdframed={style=mdblackbox}]{thmblackbox}

\declaretheoremstyle[spaceabove=6pt,spacebelow=6pt]{basehead}

\declaretheorem[style=basehead,name=Problem]{problem}
\declaretheorem[style=basehead,name=Problem,numbered=no]{problem*}
\declaretheorem[style=thmbloobox,name=Setup]{setup} % for config geo
\declaretheorem[style=thmredbox,name=Example,sibling=problem]{example}
\declaretheorem[style=thmredbox,name=$\bigstar$ Problem,sibling=problem]{niceprob}
\declaretheorem[style=thmbluebox,name=Theorem,sibling=problem]{theorem}
\declaretheorem[style=thmbluebox,name=Corollary,sibling=theorem]{corollary}
\declaretheorem[style=thmbluebox,name=Proposition,sibling=theorem]{prop}
\declaretheorem[style=thmbluebox,name=Lemma,numberwithin=problem]{lemma}
\declaretheorem[style=thmbluebox,name=Theorem,numbered=no]{theorem*}
\declaretheorem[style=thmbluebox,name=Lemma,numbered=no]{lemma*}
\declaretheorem[style=thmgreenbox,name=Claim,numberwithin=problem]{claim}
\declaretheorem[style=thmgreenbox,name=Claim,numbered=no]{claim*}
\declaretheorem[style=thmblackbox,name=Remark,numberwithin=problem]{remark}
\declaretheorem[style=thmblackbox,name=Remark,numbered=no]{remark*}
\declaretheorem[style=thmgreenbox,name=Definition,sibling=theorem]{definition}
\declaretheorem[style=thmgreenbox,name=Definition,numbered=no]{definition*}
\declaretheorem[style=thmorangebox,name=Warning,sibling=theorem]{warning}
\declaretheorem[style=thmorangebox,name=Warning,numbered=no]{warning*}
\declaretheorem[style=thmorangebox,name=Note,numbered=no]{note}
\declaretheorem[style=thmblackbox,name=Exercise,sibling=problem]{exercise}
\declaretheorem[style=thmblackbox,name=Exercise,numbered=no]{exercise*}
\declaretheorem[style=thmblackbox,name=Question,sibling=problem]{question}
\declaretheorem[style=thmblackbox,name=Question,numbered=no]{question*}

\addtolength{\textheight}{3.14cm}
\providecommand{\ol}{\overline}
\providecommand{\ul}{\underline}
\providecommand{\eps}{\varepsilon}
\providecommand{\half}{\frac{1}{2}}
\providecommand{\dang}{\measuredangle} %% Directed angle
\providecommand{\CC}{\mathbb C}
\providecommand{\FF}{\mathbb F}
\providecommand{\NN}{\mathbb N}
\providecommand{\QQ}{\mathbb Q}
\providecommand{\RR}{\mathbb R}
\providecommand{\ZZ}{\mathbb Z}
\providecommand{\dg}{^\circ}
\providecommand{\ii}{\item}
\providecommand{\setdiff}{\smallsetminus}
\providecommand{\alert}[1]{{\sffamily\textbf{\textcolor{blue}{#1}}}}
\providecommand{\inv}{^{-1}}
\providecommand{\mb}[1]{\mathbf{#1}}
\providecommand{\dif}{\;d}
\newcommand*{\horzbar}{\rule[.5ex]{2.5ex}{0.5pt}}

\reversemarginpar

\newenvironment{soln}{\begin{proof}[Solution]}{\end{proof}}
\newenvironment{walkthrough}{\noindent \textbf{Walkthrough.}}{\medskip}
\newlist{walk}{enumerate}{3}
\setlist[walk]{label=\bfseries (\alph*)}

\providecommand{\pow}{\operatorname{Pow}}
\providecommand{\vocab}[1]{{\textbf{\textcolor{ForestGreen}{#1}}}}
\providecommand{\lcm}{\operatorname{lcm}}
\providecommand{\abs}[1]{\left|#1\right|}

\usepackage{asymptote}
\begin{asydef}
	defaultpen(fontsize(10pt));
	unitsize(1cm);
	size(8cm); // set a reasonable default
	usepackage("amsmath");
	usepackage("amssymb");
	settings.tex="pdflatex";
	settings.outformat="pdf";
	// Replacement for olympiad+cse5 which is not standard
	import geometry;
	// recalibrate fill and filldraw for conics
	void filldraw(picture pic = currentpicture, conic g, pen fillpen=defaultpen, pen drawpen=defaultpen)
	{ filldraw(pic, (path) g, fillpen, drawpen); }
	void fill(picture pic = currentpicture, conic g, pen p=defaultpen)
	{ filldraw(pic, (path) g, p); }
	// some geometry
	pair foot(pair P, pair A, pair B) { return foot(triangle(A,B,P).VC); }
	pair orthocenter(pair A, pair B, pair C) { return orthocentercenter(A,B,C); }
	pair centroid(pair A, pair B, pair C) { return (A+B+C)/3; }
	// cse5 abbreviations
	path CP(pair P, pair A) { return circle(P, abs(A-P)); }
	path CR(pair P, real r) { return circle(P, r); }
	pair IP(path p, path q) { return intersectionpoints(p,q)[0]; }
	pair OP(path p, path q) { return intersectionpoints(p,q)[1]; }
	path Line(pair A, pair B, real a=0.6, real b=a) { return (a*(A-B)+A)--(b*(B-A)+B); }
	// cse5 more useful functions
	picture CC() {
		picture p=rotate(0)*currentpicture;
		currentpicture.erase();
		return p;
	}
	pair MP(Label s, pair A, pair B = plain.S, pen p = defaultpen) {
		Label L = s;
		L.s = "$"+s.s+"$";
		label(L, A, B, p);
		return A;
	}
	pair Drawing(Label s = "", pair A, pair B = plain.S, pen p = defaultpen) {
		dot(MP(s, A, B, p), p);
		return A;
	}
	path Drawing(path g, pen p = defaultpen, arrowbar ar = None) {
		draw(g, p, ar);
		return g;
	}
\end{asydef}

\usepackage{mathrsfs}

\ihead{\footnotesize MAT 362 HW07}
\ohead{\footnotesize Last compiled \today}

\title{\vspace{-2em}MAT 362 HW07}
\author{\large Aditya Pahuja}
\date{\large Due 03/27}
\begin{document}
\maketitle

\begin{problem*}3.2.1.
	Show that at a hyperbolic point,
	the principal directions bisect the asymptotic directions.
\end{problem*}

\begin{soln}
	The locus of points on the Dupin indicatrix are the points
	$ae_1 + be_2$ such that $k_1a^2 + k_2b^2 = 1$
	or $k_1a^2 + k_2b^2 = -1$
	where $e_1$, $e_2$ are the principal directions
	and $k_1$, $k_2$ the principal curvatures;
	this is a union of two hyperbolas whose (shared) asymptotes
	are the asymptotic directions at $p$.
	
	Since the axes of symmetry of the hyperbolas
	bisect their asymptotes, we just need to check that
	the axes of symmetry of $k_1a^2 + k_2b^2 = 1$ are the principal directions
	(shared asymptotes means that the axes of symmetry
	are the same for the other hyperbola).
	This is obvious, though: $e_1$ and $e_2$ are the axes of symmetry
	if and only if the maps $a\mapsto -a$ and $b\mapsto -b$
	both preserve the hyperbola, and this is trivial
	because $a^2 = (-a)^2$ and $b^2 = (-b)^2$.
\end{soln}

\begin{problem*}3.2.3.
	Let $C\subseteq S$ be a regular curve on a surface $S$
	with Gaussian curvature $K > 0$.
	Show that the curvature $k$ of $C$ at $p$ satisfies
	\[|k|\ge\min(|k_1|,|k_2|),\]
	where $k_1$ and $k_2$ are the principal curvatures of $S$ at $p$.
\end{problem*}

\begin{soln}
	Let $k_n$ be the normal curvature of $C$ at $p$.
	Since $K = k_1k_2$ is positive, $k_1$ and $k_2$ have the same sign.
	The principal curvatures are the extremal normal curvatures
	among all normal curvatures at $p$ of curves in $S$,
	so $k_n$ lies between $k_1$ and $k_2$,
	hence $k_1$ and $k_2$ having the same sign means that
	$|k_n| \ge\min(|k_1|,|k_2|)$.
	In turn, we have
	\[|k|\ge |k\cos\theta| = |k_n| \ge \min(|k_1|,|k_2|)\]
	where $\theta = \langle n,N\rangle$, where $n$ and $N$ are
	the (unit) normal vectors to $C$ and $S$ respectively at $p$.
\end{soln}

\begin{problem*}3.3.1.
	Show that at the origin $(0,0,0)$ of the hyperboloid $z=axy$
	we have $K = -a^2$ and $H = 0$.
\end{problem*}

\begin{soln}
	Let $p = (0,0,0)$ and observe that
	the hyperboloid $S$ is parametrized by $\mb x(u,v) = (u,v,auv)$.
	We have $T_p(S)$ has basis $\mb x_u = (1,0,av) = (1,0,0)$
	and $\mb x_v = (0,1,au) = (0,1,0)$.
	For a unit vector $(s,t,0)$ in $T_p(S)$,
	the plane through $(s,t,0)$ and $N(p) = (0,0,1)$ intersects $S$
	in the points $(sx, tx, y)$ satisfying $y = astx^2$,
	i.e. the curve $\alpha(x) = (sx,tx,astx^2)$. The curvature of this curve
	at $p$ is
	\[\frac{|\alpha'\wedge\alpha''|}{|\alpha'|^3}
	= \frac{|(s,t,0)\wedge(0,0,2ast)|}{|(s,t,0)|^3}
	= |(2ast^2,-2as^2t,0)| = |2ast|.\]
	The normal to this curve at $p$ is
	\[\left[\frac{d}{dx}\frac{\alpha'(x)}{|\alpha'(x)|}\right]_p
	= \frac{1}{|2ast|}
	(0,0,2ast),\]
	which is either $(0,0,1)$ or $(0,0,-1)$
	depending on if $2ast$ is positive or negative respectively,
	while the normal to $S$ at $p$ is just $(0,0,1)$.
	This means the normal curvature at $p$ is $2ast$
	(either $|2ast|\cdot 1$ if $2ast$ is positive
	or $|2ast|\cdot -1$ if $2ast$ is negative).
	Since we have pre-emptively fixed $s^2+t^2=1$ this means
	$2st$ is within the interval $[-1,1]$ because $1 = s^2+t^2 \ge |2st|$
	with equality iff $|s| = |t|$ by the AM-GM inequality,
	so the prinicpal curvatures are $\pm a$.
	Therefore, $K = -a\cdot a = \boxed{-a^2}$ and
	$H = \frac{-a + a}{2} = \boxed 0$ as desired.
\end{soln}

\pagebreak
\begin{problem*}3.3.5.
	Consider the parametrized surface
	\[\mb x(u,v) = \left(u-\frac{u^3}{3}+uv^2,
	v-\frac{v^3}{3}+vu^2,u^2-v^2\right),\]
	known as \vocab{Enneper's surface}, and show that
	\begin{enumerate}[(a)]
		\ii The coefficients of the first fundamental form are
		\[E = G = (1+u^2+v^2)^2,\quad F=0.\]
		\ii The coefficients of the second fundamental form are
		\[e=2,\quad g=-2,\quad f=0.\]
		\ii The principal curvatures are
		\[k_1 = \frac{2}{(1+u^2+v^2)^2},\quad
		k_2 = -\frac{2}{(1+u^2+v^2)^2}.\]
		\ii The lines of curvature are the coordinate curves.
		\ii The asymptotic curves are
		$u+v=\mathrm{const.}$ and $u-v=\mathrm{const.}$.
	\end{enumerate}
\end{problem*}

\begin{soln}
	(a) We have $\mb x_u = (1-u^2+v^2,2uv,2u)$
	and $\mb x_v = (2uv,1-v^2+u^2,-2v)$, so
	\begin{align*}
		E &= (1-u^2+v^2)^2 + (2uv)^2 + (2u)^2\\
		&= 1+u^4+v^4-2u^2v^2-2u^2+2v^2+4u^2v^2+4u^2\\
		&= 1+u^4+v^2+2u^2v^2+2u^2+2v^2\\
		&= (1+u^2+v^2)^2.
	\end{align*}
	The calculation for $G$ is symmetric (swap $u$ and $v$). Finally,
	\begin{align*}
		F &= 2uv(1-u^2+v^2+1+u^2-v^2) - 4uv\\
		&= 2uv\cdot 2 - 4uv = 0.
	\end{align*}

	(b) We first calculate
	\[\mb x_u\wedge\mb x_v = (-2u(u^2+v^2+1), 2v(u^2+v^2+1),1-(u^2+v^2)^2)
	= (u^2+v^2+1)\cdot(-2u,2v,1-u^2-v^2)\]
	and
	\begin{align*}
		|\mb x_u\wedge\mb x_v|
		&= (u^2+v^2+1)\sqrt{4u^2+4v^2+(u^2+v^2-1)^2}\\
		&= (u^2+v^2+1)\sqrt{u^4+v^4+1+2u^2+2v^2+2u^2v^2}\\
		&= (u^2+v^2+1)^2,
	\end{align*}
	so the normal is
	\[N = \frac{1}{u^2+v^2+1}(-2u,2v,1-u^2-v^2).\]
	Now we can compute
	\begin{align*}
		e &= \langle N,\mb x_{uu}\rangle
		= \frac{1}{u^2+v^2+1}((-2u)^2 + (2v)^2 + 2(1-u^2-v^2)) = 2\\
		g &= \langle N,\mb x_{vv}\rangle
		= \frac{1}{u^2+v^2+1}((-2u)\cdot 2u + (-2v)\cdot 2v - 2(1-u^2-v^2))
		= -2\\
		f &= \langle N,\mb x_{uv}\rangle
		= \frac{1}{u^2+v^2+1}((-2u)\cdot 2v + 2v\cdot 2u) = 0.
	\end{align*}

	(c) The principal curvatures are $H\pm\sqrt{H^2-K}$,
	which can be expanded as
	\[\half\frac{eG-2fF+gE}{EG-F^2}\pm
	\sqrt{\left(\half\frac{eG-2fF+gE}{EG-F^2}\right)^2
	- \frac{eg-f^2}{EG-F^2}}.\]
	We have $2fF = 0$ and $e=-g$, $E=G$ implying $eG+gE=0$, so
	$eG-2fF+gE=0$. Thus the principal curvatures are
	\[\pm\sqrt{-\frac{eg-f^2}{EG-F^2}}
	= \pm\sqrt{-\frac{eg}{EG}}
	= \pm\sqrt{-\frac{2\cdot(-2)}{(1+u^2+v^2)4}}
	= \frac{\pm 2}{(1+u^2+v^2)^2}.\]
	
	(d) Every point of the surface is not umbilical
	because part (c) shows that the principal curvatures are always different,
	so, since $F = f = 0$, the coordinate curves are the lines of curvature,
	as stated on page 163.
	
	(e) The asymptotic curves are the curves $\alpha(t) = \mb x(u(t), v(t))$
	satisfying the differential equation
	\[e(u')^2 + fu'v' + g(v')^2 = 2(u')^2 - 2(v')^2 =
	2(u'+v')(u'-v') = 0\]
	for all $t\in I$ (where $\alpha$ is a map $I\to\RR^3$).
	We deduce that either $u' = v'$ always holds or $u' = -v'$ always holds.
	The only alternative is that at some point either $u'$ or $v'$
	changes sign (say WLOG $u'$) while the other doesn't change sign;
	this implies that $u'(x) = 0$ for some $x\in I$,
	which then also implies $v'(x) = 0$, but then
	$\alpha'(x) = (0,0,0)$, which contradicts the regularity of $\alpha$.
	Thus $\alpha(x) = (u(x), c + v(x))$ for all $x\in I$
	or $(u(x), c - v(x))$ for all $x\in I$ where $c$ is some constant;
	either way, $\alpha$ can be reparametrized as
	one of the lines $u+v=c$ or $u-v=c$ as desired.
\end{soln}

























\end{document}